\section{Main Results}

Table~\ref{tab:main-results} now separates three questions that earlier drafts blurred together: what the strongest prompt-only single model can do, what the strongest open/local trainable model can do, and what a real hybrid workflow can do. The first point is straightforward. The stronger Chinese frontier objection is real: the DeepSeek-V3.2 few-shot dictionary baseline is a genuinely strong single-model result. It reaches 28.54 chrF++, 2.62 semantic completeness, and 2.78 fluency on the legacy repository judge, and later remains the strongest single model under the reference-aware judge.

The positive local-method result still lies inside the trainable regime. Structured multitask SFT remains the strongest pure supervised baseline. Once DPO is rerun on that base with reward-gap filtering, the story changes again: strict filtering does not just rescue DPO qualitatively, it produces the strongest local learned configuration in the repository. The clearest example is MT-DPO g0.4 sig. Tightening the reward-gap threshold from 0.2 to 0.4 cuts the preference set from 1,996 pairs to 498, raises local chrF++ from 20.93 to 30.01, and brings the average output length back near the reference ratio (1.12 instead of 1.74). The strict robust variant is slightly weaker on chrF++ but cleaner than its gap-0.2 counterpart. By contrast, the legacy DPO checkpoint remains a genuine negative control.

The strongest overall non-reference result, however, is no longer a single model. It is the catalog-aware hybrid workflow. Adjudicating among the frontier output, UNK-SFT, and MT-DPO g0.4 sig yields 34.87 chrF++, 14/50 exact matches, and the strongest legacy-judge fluency among non-reference systems in Table~\ref{tab:main-results}. This already changes the paper's practical message. The best question is not ``frontier or local?'' but ``what can local candidates add to a frontier-first workflow if adjudication is conservative and title-aware?''

\begin{table}[t]
\centering
\small
\begin{tabular}{lccccc}
\toprule
Method & Lex $\uparrow$ & PPL $\downarrow$ & chrF++ $\uparrow$ & Sem $\uparrow$ & Flu $\uparrow$ \\
\midrule
baseline1 & 0.1165 & 20.21 & 0.19 & 1.00 & 1.56 \\
baseline2 & 0.8187 & 4474.44 & 6.13 & 2.30 & 2.40 \\
baseline2\_1\_cot & 0.8723 & 18119.16 & 7.59 & 2.30 & 1.88 \\
baseline3 & 0.3192 & 298684.91 & 17.85 & 1.42 & 1.88 \\
baseline3\_1\_unk & 0.3514 & 11031.18 & 22.11 & 1.66 & 1.96 \\
baseline3\_2\_multitask & 0.4839 & 3322.93 & \textbf{25.99} & 1.70 & 2.26 \\
baseline3\_3\_semantic & 0.3723 & 17064.38 & 17.03 & 1.46 & 1.84 \\
final\_v2 & 0.4224 & 314577.16 & 19.39 & 1.52 & 1.76 \\
final\_gap02\_multitask\_sigmoid & \textbf{0.5694} & 303438.00 & 20.93 & 1.88 & 1.86 \\
final\_gap02\_multitask\_robustwpo & 0.5249 & 238892.42 & \textbf{27.67} & 1.88 & 1.88 \\
human\_reference & 0.5733 & 1646139.34 & 100.00 & 2.24 & 2.34 \\
\bottomrule
\end{tabular}
\caption{Repository-level metrics on the 50-example Tangut test set. The strongest pure SFT baseline is \texttt{baseline3\_2\_multitask}, while multitask-base DPO exposes a metric trade-off: \texttt{final\_gap02\_multitask\_robustwpo} gives the highest raw chrF++, but later reference-aware judging prefers \texttt{final\_gap02\_multitask\_sigmoid}.}
\label{tab:main-results}
\end{table}


Not all of that workflow story depends on a closed adjudicator. A deterministic open selector over the same frontier/UNK-SFT/MT-DPO pool, described in the methods section but omitted from Table~\ref{tab:main-results} because it does not use the legacy repository submetrics as primary evidence, still reaches 32.62 chrF++, 12/50 exact matches, and zero contamination. This does not match the full catalog hybrid, but it shows that part of the gain is reproducible with a cheap open reranker rather than only with GPT-based adjudication.

These disagreements are precisely why the paper should not rely on a single automatic metric. Raw chrF++, contamination, and the reference-aware judge each reward different parts of the short-title problem. The defensible claim is therefore two-level rather than absolutist: strict gap-filtered multitask DPO is still the best local training recipe, but catalog-aware adjudication turns that local complementarity into the strongest overall workflow result.

The new statistical analysis makes both layers more precise. Under 5,000 paired bootstrap resamples, the best local DPO model improves over multitask SFT by +4.02 chrF++ with a 95\% confidence interval of [0.84, 7.80], so the local chrF++ gain is not just a single-split fluctuation. More importantly for the revised paper, the catalog hybrid improves over the frontier baseline by +6.33 chrF++ with a 95\% confidence interval of [2.10, 11.67]. The gains in reference-aware overall (+0.22, 95\% CI [0.00, 0.48]) and exact match (+0.04, 95\% CI [0.00, 0.10]) are positive but boundary-level under this 50-example test set. This is still a strong workflow result: the hybrid is clearly better on raw overlap, clearly no worse on the reference-aware judge, and directionally better on exact match, while keeping zero contamination and near-reference title length.
